\documentclass[10pt,a4paper]{amsart}
\usepackage{amsmath,amssymb,amsthm,graphicx}
% Retained preamble for the Stacks Project chapter divisors.tex (Divisors).
% The chapter file pinned-original-source.tex carries only \input{preamble} on its first line, so no
% preamble text lives inside the chapter itself.  The Stacks Project uses one single shared preamble for
% every chapter; the snapshot stages it verbatim in sources/stacks-chow/preamble.tex.  The block below is
% copied mechanically, line for line, from that file's lines 14-18 (the Xy-pic setup) and lines 154-217
% (the native theorem declarations together with the mathematical macros).  Nothing was added, removed or
% redefined.  The wrapper class and the amsmath/amssymb/amsthm/graphicx base packages are supplied by the
% assembler, so the class/\externaldocument part of the shared preamble is not repeated here.
% For commutative diagrams we use Xy-pic
\usepackage[all]{xy}

% We use 2cell for 2-commutative diagrams.
\xyoption{2cell}
% Theorem environments.
%
\theoremstyle{plain}
\newtheorem{theorem}[subsection]{Theorem}
\newtheorem{proposition}[subsection]{Proposition}
\newtheorem{lemma}[subsection]{Lemma}

\theoremstyle{definition}
\newtheorem{definition}[subsection]{Definition}
\newtheorem{example}[subsection]{Example}
\newtheorem{exercise}[subsection]{Exercise}
\newtheorem{situation}[subsection]{Situation}

\theoremstyle{remark}
\newtheorem{remark}[subsection]{Remark}
\newtheorem{remarks}[subsection]{Remarks}

\numberwithin{equation}{subsection}

% Macros
%
\def\lim{\mathop{\mathrm{lim}}\nolimits}
\def\colim{\mathop{\mathrm{colim}}\nolimits}
\def\Spec{\mathop{\mathrm{Spec}}}
\def\Hom{\mathop{\mathrm{Hom}}\nolimits}
\def\Ext{\mathop{\mathrm{Ext}}\nolimits}
\def\SheafHom{\mathop{\mathcal{H}\!\mathit{om}}\nolimits}
\def\SheafExt{\mathop{\mathcal{E}\!\mathit{xt}}\nolimits}
\def\Sch{\mathit{Sch}}
\def\Mor{\mathop{\mathrm{Mor}}\nolimits}
\def\Ob{\mathop{\mathrm{Ob}}\nolimits}
\def\Sh{\mathop{\mathit{Sh}}\nolimits}
\def\NL{\mathop{N\!L}\nolimits}
\def\CH{\mathop{\mathrm{CH}}\nolimits}
\def\proetale{{pro\text{-}\acute{e}tale}}
\def\etale{{\acute{e}tale}}
\def\QCoh{\mathit{QCoh}}
\def\Ker{\mathop{\mathrm{Ker}}}
\def\Im{\mathop{\mathrm{Im}}}
\def\Coker{\mathop{\mathrm{Coker}}}
\def\Coim{\mathop{\mathrm{Coim}}}

% Boxtimes
%
\DeclareMathSymbol{\boxtimes}{\mathbin}{AMSa}{"02}

%
% Macros for moduli stacks/spaces
%
\def\QCohstack{\mathcal{QC}\!\mathit{oh}}
\def\Cohstack{\mathcal{C}\!\mathit{oh}}
\def\Spacesstack{\mathcal{S}\!\mathit{paces}}
\def\Quotfunctor{\mathrm{Quot}}
\def\Hilbfunctor{\mathrm{Hilb}}
\def\Curvesstack{\mathcal{C}\!\mathit{urves}}
\def\Polarizedstack{\mathcal{P}\!\mathit{olarized}}
\def\Complexesstack{\mathcal{C}\!\mathit{omplexes}}
% \Pic is the operator that assigns to X its picard group, usage \Pic(X)
% \Picardstack_{X/B} denotes the Picard stack of X over B
% \Picardfunctor_{X/B} denotes the Picard functor of X over B
\def\Pic{\mathop{\mathrm{Pic}}\nolimits}
\def\Picardstack{\mathcal{P}\!\mathit{ic}}
\def\Picardfunctor{\mathrm{Pic}}
\def\Deformationcategory{\mathcal{D}\!\mathit{ef}}
\usepackage{cleveref}
\crefname{originalreference}{original reference}{original references}
\Crefname{originalreference}{Original reference}{Original references}
\makeatletter
\newlabel{algebra-lemma-Noetherian-irreducible-components}{{algebra-lemma-Noetherian-irreducible-components}{source}{}{sourcepointer-1}{}}
\newlabel{algebra-lemma-Noetherian-irreducible-components@cref}{{[originalreference][0][]algebra-lemma-Noetherian-irreducible-components}{[1][0][]source}{}{}{}}
\makeatother
\crefname{originalreference}{original reference}{original references}
\Crefname{originalreference}{Original reference}{Original references}
\makeatletter
\newlabel{algebra-lemma-going-up-closed}{{algebra-lemma-going-up-closed}{source}{}{sourcepointer-2}{}}
\newlabel{algebra-lemma-going-up-closed@cref}{{[originalreference][0][]algebra-lemma-going-up-closed}{[1][0][]source}{}{}{}}
\makeatother
\crefname{originalreference}{original reference}{original references}
\Crefname{originalreference}{Original reference}{Original references}
\makeatletter
\newlabel{algebra-lemma-integral-going-up}{{algebra-lemma-integral-going-up}{source}{}{sourcepointer-3}{}}
\newlabel{algebra-lemma-integral-going-up@cref}{{[originalreference][0][]algebra-lemma-integral-going-up}{[1][0][]source}{}{}{}}
\makeatother
\crefname{originalreference}{original reference}{original references}
\Crefname{originalreference}{Original reference}{Original references}
\makeatletter
\newlabel{algebra-lemma-one-equation}{{algebra-lemma-one-equation}{source}{}{sourcepointer-4}{}}
\newlabel{algebra-lemma-one-equation@cref}{{[originalreference][0][]algebra-lemma-one-equation}{[1][0][]source}{}{}{}}
\makeatother
\crefname{originalreference}{original reference}{original references}
\Crefname{originalreference}{Original reference}{Original references}
\makeatletter
\newlabel{algebra-lemma-silly}{{algebra-lemma-silly}{source}{}{sourcepointer-5}{}}
\newlabel{algebra-lemma-silly@cref}{{[originalreference][0][]algebra-lemma-silly}{[1][0][]source}{}{}{}}
\makeatother
\crefname{originalreference}{original reference}{original references}
\Crefname{originalreference}{Original reference}{Original references}
\makeatletter
\newlabel{algebra-lemma-support-point}{{algebra-lemma-support-point}{source}{}{sourcepointer-6}{}}
\newlabel{algebra-lemma-support-point@cref}{{[originalreference][0][]algebra-lemma-support-point}{[1][0][]source}{}{}{}}
\makeatother
\title{The punctured spectrum of a Noetherian local ring is affine if and only if the dimension is at most one}
\author{The Stacks Project Authors}
\date{stacks-project:divisors.tex:local-snapshot-20261012; GNU FDL 1.2 or later}
\begin{document}
\maketitle
\noindent\textit{Editorial note.} All delivered mathematical text is the author's own, copied from the pinned Stacks Project chapter file pinned-original-source.tex; no statement or proof step was rewritten, reordered or removed.\par\medskip
\noindent\textit{Editorial note.} Following the multi-claim convention for this source, source\_url is the source-level URL https://stacks.math.columbia.edu/ and the source label lemma-affine-punctured-spec is carried in claim\_key and here rather than in a per-claim tag URL.\par\medskip
\noindent\textit{Editorial note.} The preamble is the chapter's shared Stacks Project preamble, retained as sources/stacks-divisors/retained-preamble.tex: a line-for-line excerpt of the shared Stacks preamble staged with this snapshot. No macro was invented or redefined, and this card adds no preamble text of its own.\par\medskip
\noindent\textit{Editorial note.} Statement span 3348-3353 and proof span 3355-3407 come from the halfA pool card operations/pool-cards/stacks-divisors-halfA.json and were re-read in the pinned chapter file before this card was written.\par\medskip
\noindent\textit{Editorial note.} The author announces in the proof that the argument given is 'somewhat nonstandard' and that the standard argument is discussed in a remark below; that sentence is kept verbatim, and the remark it refers to is outside the counted span and is not part of this card.\par\medskip
\noindent\textit{Editorial note.} The proof writes the snake lemma sequence and the two length sequences with the author's own notation (Q[f], l\_n); all displays were copied unchanged.\par\medskip
\noindent\textit{Editorial note.} 6 cross-references occur inside the counted spans; they are listed in external\_references as navigation pointers only.\par\medskip
\noindent\textit{Allowed standard prerequisites (editorial).}\begin{itemize}
\item A Noetherian ring has finitely many minimal primes, and an element of the maximal ideal avoiding them can be chosen; the spectrum of a localisation is the corresponding open subset. Cited in the same chapter's algebra parent at proof lines 3358-3361 and reused at 3377-3379. Produces the element f that makes the punctured spectrum a principal affine open in dimension one, and the nonzerodivisor used in the converse.
\item Dimension drops by one on cutting out a nonzerodivisor, and a finite ring map induces a closed map on spectra. Cited at proof lines 3382-3383 and again at 3398-3400. Gives the dimension count for A modulo fA and turns the finiteness of the ring map A/fA to B/fB into closedness of the induced map on spectra, which produces the contradiction.
\item A module whose support is the closed point of a Noetherian local ring has finite length. Cited at proof lines 3390-3391. Makes the length l of the module of f-torsion in Q finite, which is what bounds the lengths of the modules Q[f\textasciicircum{}n].
\end{itemize}\medskip
\ifdefined\hypertarget\hypertarget{sourcepointer-1}{}\fi
\noindent\textit{Source reference.} Original reference algebra-lemma-Noetherian-irreducible-components: sources/stacks-divisors/pinned-original-source.tex, source lines 3360--3360. Same-chapter reference used in the retained author text at source line(s) 3360. It is supplied only as a navigation pointer so that the retained text compiles; the counted item delivers the author's own statement and the author's complete proof, and no part of the referenced item is copied in or substituted for the counted argument.\par\medskip
\ifdefined\hypertarget\hypertarget{sourcepointer-2}{}\fi
\noindent\textit{Source reference.} Original reference algebra-lemma-going-up-closed: sources/stacks-divisors/pinned-original-source.tex, source lines 3403--3403. Same-chapter reference used in the retained author text at source line(s) 3403. It is supplied only as a navigation pointer so that the retained text compiles; the counted item delivers the author's own statement and the author's complete proof, and no part of the referenced item is copied in or substituted for the counted argument.\par\medskip
\ifdefined\hypertarget\hypertarget{sourcepointer-3}{}\fi
\noindent\textit{Source reference.} Original reference algebra-lemma-integral-going-up: sources/stacks-divisors/pinned-original-source.tex, source lines 3402--3402. Same-chapter reference used in the retained author text at source line(s) 3402. It is supplied only as a navigation pointer so that the retained text compiles; the counted item delivers the author's own statement and the author's complete proof, and no part of the referenced item is copied in or substituted for the counted argument.\par\medskip
\ifdefined\hypertarget\hypertarget{sourcepointer-4}{}\fi
\noindent\textit{Source reference.} Original reference algebra-lemma-one-equation: sources/stacks-divisors/pinned-original-source.tex, source lines 3036--3036. Same-chapter reference used in the retained author text at source line(s) 3036, 3378, 5578. It is supplied only as a navigation pointer so that the retained text compiles; the counted item delivers the author's own statement and the author's complete proof, and no part of the referenced item is copied in or substituted for the counted argument.\par\medskip
\ifdefined\hypertarget\hypertarget{sourcepointer-5}{}\fi
\noindent\textit{Source reference.} Original reference algebra-lemma-silly: sources/stacks-divisors/pinned-original-source.tex, source lines 3361--3361. Same-chapter reference used in the retained author text at source line(s) 3361, 6306. It is supplied only as a navigation pointer so that the retained text compiles; the counted item delivers the author's own statement and the author's complete proof, and no part of the referenced item is copied in or substituted for the counted argument.\par\medskip
\ifdefined\hypertarget\hypertarget{sourcepointer-6}{}\fi
\noindent\textit{Source reference.} Original reference algebra-lemma-support-point: sources/stacks-divisors/pinned-original-source.tex, source lines 3387--3387. Same-chapter reference used in the retained author text at source line(s) 3387. It is supplied only as a navigation pointer so that the retained text compiles; the counted item delivers the author's own statement and the author's complete proof, and no part of the referenced item is copied in or substituted for the counted argument.\par\medskip
\begin{lemma}
\label{lemma-affine-punctured-spec}
Let $(A, \mathfrak m)$ be a Noetherian local ring.
The punctured spectrum $U = \Spec(A) \setminus \{\mathfrak m\}$
of $A$ is affine if and only if $\dim(A) \leq 1$.
\end{lemma}

\begin{proof}
If $\dim(A) = 0$, then $U$ is empty hence affine (equal to the spectrum of
the $0$ ring). If $\dim(A) = 1$, then we can choose an element
$f \in \mathfrak m$ not contained in any of the finite number of minimal
primes of $A$
(Algebra, Lemmas \ref{algebra-lemma-Noetherian-irreducible-components} and
\ref{algebra-lemma-silly}). Then $U = \Spec(A_f)$
is affine.

\medskip\noindent
The converse is more interesting. We will give a somewhat nonstandard proof
and discuss the standard argument in a remark below.
Assume $U = \Spec(B)$ is affine. Since affineness and dimension are not
affected by going to the reduction we may replace $A$ by the quotient by
its ideal of nilpotent elements and assume $A$ is reduced.
Set $Q = B/A$ viewed as an $A$-module.
The support of $Q$ is $\{\mathfrak m\}$ as $A_\mathfrak p = B_\mathfrak p$
for all nonmaximal primes $\mathfrak p$ of $A$.
We may assume $\dim(A) \geq 1$, hence as above we can pick
$f \in \mathfrak m$ not contained in any of the minimal ideals of $A$.
Since $A$ is reduced this implies that $f$ is a nonzerodivisor in $A$.
The same argument shows that $f : B \to B$ is injective as $\Spec(B)$
is open in $\Spec(A)$. Note that $\dim(A/fA) = \dim(A) - 1$ by
Algebra, Lemma \ref{algebra-lemma-one-equation}.
Applying the snake lemma to multiplication by $f$ on the short
exact sequence $0 \to A \to B \to Q \to 0$ we obtain
$$
0 \to Q[f] \to A/fA \to B/fB \to Q/fQ \to 0
$$
where $Q[f] = \Ker(f : Q \to Q)$.
This implies that $Q[f]$ is a finite $A$-module. Since the support of
$Q[f]$ is $\{\mathfrak m\}$ we see $l = \text{length}_A(Q[f]) < \infty$
(Algebra, Lemma \ref{algebra-lemma-support-point}).
Set $l_n = \text{length}_A(Q[f^n])$. The exact sequence
$$
0 \to Q[f^n] \to Q[f^{n + 1}] \xrightarrow{f^n} Q[f]
$$
shows inductively that $l_n < \infty$ and that $l_n \leq l_{n + 1}$.
Considering the exact sequence
$$
0 \to Q[f] \to Q[f^{n + 1}] \xrightarrow{f} Q[f^n] \to Q/fQ
$$
and we see that the image of $Q[f^n]$ in $Q/fQ$ has length
$l_n - l_{n + 1} + l \leq l$. Since $Q = \bigcup Q[f^n]$ we
find that the length of $Q/fQ$ is at most $l$, i.e., bounded.
Thus $Q/fQ$ is a finite $A$-module. Hence $A/fA \to B/fB$ is a
finite ring map, in particular induces a closed map on spectra
(Algebra, Lemmas \ref{algebra-lemma-integral-going-up} and
\ref{algebra-lemma-going-up-closed}).
On the other hand $\Spec(B/fB)$ is the punctured spectrum of $\Spec(A/fA)$.
This is a contradiction unless $\Spec(B/fB) = \emptyset$ which
means that $\dim(A/fA) = 0$ as desired.
\end{proof}

\end{document}
